\documentclass[12pt,letterpaper]{article}

% ── Packages ─────────────────────────────────────────────────────────────────
\usepackage{amsmath,amssymb,amsthm}
\usepackage{mathtools}
\usepackage{geometry}
\usepackage{hyperref}
\usepackage{booktabs}
\usepackage{array}
\usepackage{graphicx}
\usepackage{xcolor}
\usepackage{microtype}
\usepackage[numbers,sort&compress]{natbib}

\geometry{margin=1in}

% ── Theorem environments ──────────────────────────────────────────────────────
\theoremstyle{plain}
\newtheorem{theorem}{Theorem}[section]
\newtheorem{lemma}[theorem]{Lemma}
\newtheorem{proposition}[theorem]{Proposition}
\newtheorem{corollary}[theorem]{Corollary}

\theoremstyle{definition}
\newtheorem{definition}[theorem]{Definition}
\newtheorem{remark}[theorem]{Remark}

% ── Custom commands ───────────────────────────────────────────────────────────
\newcommand{\Hilb}{\mathcal{H}}
\newcommand{\Dens}{\mathcal{D}(\mathcal{H})}
\newcommand{\Cman}{\mathcal{C}}
\newcommand{\Obs}{\mathcal{O}}
\newcommand{\tr}{\operatorname{Tr}}
\newcommand{\SvN}{S_{\mathrm{vN}}}
\newcommand{\FQ}{F_{Q}}
\newcommand{\Bspam}{B_{\mathrm{SPAM}}}
\newcommand{\hSvN}{\hat{S}_{\mathrm{vN}}}
\newcommand{\half}{\tfrac{1}{2}}

% ── Metadata ─────────────────────────────────────────────────────────────────
\title{\textbf{Constraint Manifolds and the Limits of Quantum Observability}\\[4pt]
\large Resolving the 40--50\% Entropy Plateau in NISQ Devices}

\author{Kevin Monette\\[2pt]
\small Independent Researcher (AI-assisted research)\\
\small \texttt{kevin.monette@research.org}}

\date{February 9, 2026 \quad (repository version March 2026)}

% ─────────────────────────────────────────────────────────────────────────────
\begin{document}
\maketitle

\begin{abstract}
In near-term quantum devices with 16--28$+$ qubits, state tomography consistently yields
von~Neumann entropy estimates at 40--50\% of the theoretical maximum --- a plateau widely
attributed to a fundamental decoherence limit of current hardware.  We show instead that
this plateau is a \emph{measurement artifact arising from finite-sample estimator bias}, not a
physical decoherence ceiling.

We formalize the distinction through the constraint manifold
$\Cman \subset \Dens$ --- the subspace of physically accessible density operators --- and
prove that required measurement resources scale as $\nu \propto 2^{n/2}$ via quantum Fisher
information analysis.  For typical NISQ tomography budgets ($\nu \sim 10^{3}$ shots) on
systems with $n \geq 15$ qubits, finite-sample estimator bias alone accounts for the observed
plateau without invoking physical decoherence.

The framework connects naturally to thermodynamic gravity: constraint-induced entropy
gradients source effective stress-energy via the coupling derived in a companion
paper~\cite{MonetteA2026}.

\medskip\noindent
\textbf{Falsification criterion.} If entropy estimates for $n$-qubit coherent states fail to
converge to $\SvN < 0.1$\,bits when $\nu \geq 100 \cdot 2^{n/2}$ after SPAM calibration, the
measurement-insufficiency hypothesis is falsified.
\end{abstract}

\tableofcontents
\bigskip

% ─────────────────────────────────────────────────────────────────────────────
\section{The Constraint Manifold Formalism}
\label{sec:manifold}

\subsection{Physically Accessible State Space}

Not all density operators in $\Dens$ are physically accessible to a quantum system evolving
under real conditions.  Conservation laws, decoherence channels, and initial-state constraints
restrict evolution to a proper sub-manifold.

\begin{definition}[Constraint Manifold]
\label{def:manifold}
\begin{equation}
  \Cman
  \;=\;
  \bigl\{\,\rho \in \Dens
      \;\big|\;
      \tr\!\bigl(\hat{C}_{i}\,\rho\bigr) = c_{i}
      \;\;\forall\, i \in \mathcal{I}_{\mathrm{irr}}
  \bigr\},
  \label{eq:constraint-manifold}
\end{equation}
where $\hat{C}_{i}$ are constraint operators (e.g.\ $\hat{H}$ for energy conservation,
$\hat{Q}$ for charge), $c_{i}$ are the constraint values fixed by initial conditions, and
$\mathcal{I}_{\mathrm{irr}}$ indexes \emph{irreversible} constraints --- those that cannot be
undone by subsequent unitary evolution.
\end{definition}

Soft constraints (e.g.\ thermodynamic bias toward equilibrium) enter through a measure $\mu$
on $\Cman$ rather than its definition:
\begin{equation}
  \mu(d\rho)
  \;\propto\;
  e^{-\beta\,\tr(\hat{H}\rho)}\,d\rho.
  \label{eq:soft-constraint}
\end{equation}

\subsection{The Observable Subspace}

A measurement history $r = \{i_{1}, i_{2}, \ldots, i_{k}\}$ represents the sequence of
constraint-fixing events that have become irreversible.  Each event reduces the accessible
subspace:
\begin{equation}
  \Obs(r)
  \;=\;
  \bigl\{\,\rho \in \Cman
      \;\big|\;
      \tr\!\bigl(\hat{C}_{i_{j}}\,\rho\bigr) = c_{i_{j}}
      \;\;\forall\, j \leq k
  \bigr\}.
  \label{eq:observable-subspace}
\end{equation}

As $k$ increases, $\Obs(r)$ shrinks.  This is not decoherence --- it is the progressive fixing
of constraints by observation.

\begin{proposition}[Dimension of the Observable Subspace]
\label{prop:dimension}
For an $n$-qubit system with $m$ linearly independent measurement constraints,
\begin{equation}
  \dim\Obs(r)
  \;=\;
  4^{n} - \operatorname{rank}(\mathcal{C}) - 1,
  \label{eq:dim}
\end{equation}
where $\operatorname{rank}(\mathcal{C})$ is the number of linearly independent constraints.
\end{proposition}

\begin{proof}
The space of $n$-qubit density operators (Hermitian, positive semidefinite, trace-1) has
real dimension $4^{n}-1$ (the $2^{n}\times 2^{n}$ Hermitian matrix parameters minus the
trace normalization).  Each linearly independent irreversible constraint reduces the dimension
by one.  Linear independence must be verified in the space of Hermitian operators; redundant
constraints reduce the effective rank below $m$.
\end{proof}

\begin{remark}
When $m \ll 4^{n}$, the observable subspace vastly undersamples the physical state space,
creating an apparent entropy increase without any physical decoherence.
\end{remark}

% ─────────────────────────────────────────────────────────────────────────────
\section{Quantum Fisher Information and Shot Scaling}
\label{sec:qfi}

\subsection{The Quantum Cram\'{e}r--Rao Bound}

The quantum Cram\'{e}r--Rao bound places a fundamental lower limit on the variance of any
unbiased estimator $\hSvN$ of the von~Neumann entropy:
\begin{equation}
  \operatorname{Var}[\hSvN]
  \;\geq\;
  \frac{\bigl[\FQ(\rho)\bigr]^{-1}}{\nu},
  \label{eq:cramer-rao}
\end{equation}
where $\nu$ is the number of independent measurement shots and
$\FQ(\rho) = \tr[\rho\,L^{2}]$ is the quantum Fisher information, with $L$ the symmetric
logarithmic derivative satisfying $\partial_{\theta}\rho = (\rho L + L\rho)/2$.

\subsection{Exponential Scaling Near Maximum Entropy}

\begin{lemma}[Fisher Information Scaling]
\label{lem:fisher}
For states near the maximally mixed state $\rho \approx \mathbf{I}/2^{n}$,
\begin{equation}
  \FQ(\rho) \;\sim\; 2^{-n/2}.
  \label{eq:fisher-scaling}
\end{equation}
\end{lemma}

Near maximum entropy, all probability distributions over measurement outcomes are nearly
identical.  Distinguishing $\SvN = 0.99\,S_{\max}$ from $\SvN = S_{\max}$ requires resolving
tiny differences in an exponentially large space of nearly equally likely outcomes.  The
sensitivity of the measurement record to entropy changes scales as the average overlap between
adjacent density matrices in the entropy-variation direction, which is exponentially close to
unity in any product measurement basis.

\subsection{Required Shot Count}

To achieve precision $\varepsilon$ in the entropy estimate, the Cram\'{e}r--Rao bound requires:
\begin{equation}
  \nu
  \;\gtrsim\;
  \varepsilon^{-2}\cdot 2^{n/2}.
  \label{eq:shot-scaling}
\end{equation}

\begin{table}[h]
\centering
\caption{Required measurement shots for entropy estimation at precision $\varepsilon = 0.1$\,bits.}
\label{tab:shots}
\begin{tabular}{@{}ccccc@{}}
\toprule
Qubit count $n$ & Hilbert dim.\ $d=2^{n}$ & Required $\nu$ & Typical NISQ $\nu$ & Gap factor \\
\midrule
5  & 32          & $\approx$560       & $10^{3}$ & $<$1 (sufficient) \\
10 & 1\,024      & $\approx$3\,160    & $10^{3}$ & $\approx$3$\times$ \\
15 & 32\,768     & $\approx$17\,800   & $10^{3}$ & $\approx$18$\times$ \\
20 & 1\,048\,576 & $\approx$100\,000  & $10^{3}$ & $\approx$100$\times$ \\
25 & $\approx33$M & $\approx$560\,000 & $10^{3}$ & $\approx$560$\times$ \\
28 & $\approx268$M & $\approx$1.6M    & $10^{3}$ & $\approx$1\,600$\times$ \\
\bottomrule
\end{tabular}
\end{table}

The gap between required and typical shot budgets explains the entropy plateau for
$n \geq 15$ entirely: NISQ devices are being measured with 18--1\,600$\times$ too few shots to
accurately estimate entropy.

% ─────────────────────────────────────────────────────────────────────────────
\section{The Estimator Bias Formula}
\label{sec:bias}

\subsection{Linear Inversion Bias}

The standard linear-inversion tomographic estimator exhibits systematic bias scaling with
Hilbert space dimension $d = 2^{n}$~\cite{Flammia2012}:
\begin{equation}
  \mathbb{E}[\hSvN]
  \;=\;
  \SvN(\rho)
  \;+\;
  \underbrace{\frac{d-1}{2\nu}}_{\text{finite-sampling bias}}
  \;+\;
  \underbrace{\Bspam}_{\text{readout errors}}
  \;+\;
  O\!\left(\nu^{-2}\right).
  \label{eq:bias}
\end{equation}

The estimator \emph{systematically overestimates} entropy by $(d-1)/(2\nu)$ bits, independent
of the true entropy and regardless of hardware quality.

\subsection{Quantitative Explanation of the 40--50\% Plateau}

\paragraph{Example: $n = 15$ qubits, $\nu = 10^{4}$ shots.}
\begin{align}
  d &= 2^{15} = 32\,768, \notag\\[2pt]
  \text{Finite-sampling bias} &= \frac{d-1}{2\nu}
      = \frac{32\,767}{20\,000} \approx 1.64\;\text{bits}, \notag\\[2pt]
  S_{\max} &= n\ln 2 \approx 10.4\;\text{bits}, \notag\\[2pt]
  \text{Bias as fraction of }S_{\max} &\approx 15.8\%.
\end{align}

Adding SPAM errors ($\Bspam \approx 0.5$--$1.0$\,bits for current hardware):
\begin{equation}
  \text{Total apparent inflation} \approx
  \frac{1.64 + 0.75}{10.4} \approx 23\%.
\end{equation}
For a state with true $\SvN \approx 20$--$25\%$ of $S_{\max}$, the measured value lands at
$\approx 43$--$48\%$ --- precisely matching the observed NISQ plateau without physical
decoherence.

\subsection{The Bias-Corrected Estimator}

Define the corrected estimator:
\begin{equation}
  \hSvN^{\,\mathrm{corr}}
  \;=\;
  \hSvN - \frac{d-1}{2\nu} - \hat{B}_{\mathrm{SPAM}},
  \label{eq:corrected}
\end{equation}
where $\hat{B}_{\mathrm{SPAM}}$ is estimated from measurement calibration circuits.

Under the measurement-insufficiency hypothesis, the corrected estimator satisfies:
\begin{equation}
  \hSvN^{\,\mathrm{corr}}(\nu)
  \;=\;
  \SvN^{\,\mathrm{phys}}(\rho)
  \;+\;
  \frac{d-1}{2\nu}
  \;+\;
  O\!\left(\nu^{-2}\right),
  \label{eq:convergence}
\end{equation}
converging monotonically to the true physical entropy as $\nu \to \infty$.

% ─────────────────────────────────────────────────────────────────────────────
\section{Connection to Thermodynamic Gravity}
\label{sec:gravity}

The constraint manifold formalism connects naturally to the entropic gravity program.  In
Jacobson's thermodynamic derivation~\cite{Jacobson1995}, spacetime geometry emerges from
entropy gradients across causal horizons.  The companion paper~\cite{MonetteA2026} extends
this to laboratory scales.

Constraint-induced entropy gradients $\nabla\SvN$ source effective stress-energy via:
\begin{equation}
  \text{(entropy gradient)}
  \;\longrightarrow\;
  \text{effective stress-energy}
  \;\longrightarrow\;
  \text{spacetime curvature}.
\end{equation}

For the gravitational prediction to be testable, experiments must measure \emph{physical}
entropy changes, not artifacts of the measurement budget.  Accordingly, the shot-scaling
requirement~\eqref{eq:shot-scaling} is a prerequisite for any entanglement-gravity experiment:
SPAM-corrected, high-$\nu$ tomography must be performed before interpreting any entropy
difference as a gravitational source.

\paragraph{Observation is not passive.}
An observation is any physical process that irreversibly correlates with a constraint value of
the observed system.  Decoherence is environmental observation.  Conscious measurement is the
same physical process, more controlled.  Consciousness plays no special role --- the
constraint-fixing is thermodynamic, not mental.

% ─────────────────────────────────────────────────────────────────────────────
\section{Experimental Protocol for Falsification}
\label{sec:experiment}

\subsection{Three-Stage Validation Protocol}

\paragraph{Stage 1 --- Calibration.}
Prepare known-entropy states (computational basis states, Bell pairs, product states) with
theoretically exact $\SvN$.  Construct the SPAM correction matrix $\Lambda$ from calibration
circuits.  Verify that $\hSvN^{\,\mathrm{corr}}$ converges to theoretical $\SvN$ at rate
$O(1/\nu)$.

\paragraph{Stage 2 --- Scaling test.}
Prepare $n$-qubit GHZ states
$|\mathrm{GHZ}\rangle = (|0\rangle^{\otimes n} + |1\rangle^{\otimes n})/\!\sqrt{2}$
for $n \in \{5, 8, 10, 12, 15\}$.  Measure
$\hSvN(\nu)$ for $\nu \in \{10^{3}, 10^{4}, 10^{5}, 10^{6}\}$ shots.
Apply SPAM correction: $\hat{\rho}_{\mathrm{corr}} = \Lambda^{-1}\hat{\rho}_{\mathrm{raw}}$.

\paragraph{Stage 3 --- Convergence verification.}
Compute $\hSvN^{\,\mathrm{corr}}(\nu)$ from~\eqref{eq:corrected}.  Fit to the scaling model:
\begin{equation}
  \hSvN^{\,\mathrm{corr}}(\nu)
  = \SvN^{\,\mathrm{phys}}
  + \frac{A}{\nu}
  + O\!\left(\nu^{-2}\right).
  \label{eq:scaling-model}
\end{equation}
Extract $\SvN^{\,\mathrm{phys}}$ as the $\nu \to \infty$ extrapolation and compare to the
theoretical prediction ($\SvN = 0$ for pure GHZ).

\subsection{Falsification Criterion}

\begin{center}
\fbox{\parbox{0.88\textwidth}{%
\textbf{Falsification:} The measurement-insufficiency hypothesis is falsified if, for
$n$-qubit GHZ states with verified entanglement depth $\geq n$, the SPAM-corrected entropy
estimate $\hSvN^{\,\mathrm{corr}}$ fails to converge to $\SvN < 0.1$\,bits when shot count
$\nu \geq 100 \cdot 2^{n/2}$, verified by bootstrap resampling with $>95\%$ confidence.%
}}
\end{center}

% ─────────────────────────────────────────────────────────────────────────────
\section{Implications}
\label{sec:implications}

\subsection{For Quantum Computing}
Many NISQ devices have better coherence properties than their entropy estimates suggest.
Quantum-advantage assessments based on entropy metrics should account for the estimator bias
of Eq.~\eqref{eq:bias} before concluding that hardware is at a decoherence limit.

\subsection{For Quantum Gravity Experiments}
Physical entropy measurement (not estimated entropy) is a prerequisite for any entanglement-gravity experiment.  The shot-scaling requirement~\eqref{eq:shot-scaling} must be satisfied
before any differential entropy is attributed to a gravitational source.

\subsection{For Quantum Information}
The scaling law $\nu \propto 2^{n/2}$ is a fundamental resource bound for entropy estimation
in $n$-qubit systems.  This is sub-exponential compared to full state tomography ($\nu \propto
4^{n}$) and can in principle be achieved by randomized measurement
protocols~\cite{Cotler2020}.

% ─────────────────────────────────────────────────────────────────────────────
\section{Conclusion}
\label{sec:conclusion}

The apparent entropy plateau at 40--50\% in NISQ devices is a consequence of finite-sample
estimator bias, not fundamental hardware limitations.  Required shot count for accurate entropy
estimation scales as $\nu \propto 2^{n/2}$, derived rigorously from quantum Fisher information
analysis.  The measurement-insufficiency hypothesis makes precise, falsifiable predictions
verifiable with existing hardware.

The framework establishes three results for quantum physics broadly:
\begin{enumerate}
  \item Physical entropy and measured entropy are distinct quantities, separated by a
        computable, correctable bias.
  \item Observation is a physical resource (measurement shots) that must scale with Hilbert
        space dimension.
  \item Much apparent disorder in NISQ systems is epistemic (measurement-limited), not
        physical (decoherence-limited).
\end{enumerate}

% ─────────────────────────────────────────────────────────────────────────────
\section*{Acknowledgments}
The author acknowledges AI research partners for assistance with literature synthesis, equation
formatting, and critical review.  All theoretical content and experimental design originate from
the author's independent research program.

% ─────────────────────────────────────────────────────────────────────────────
\bibliographystyle{unsrtnat}
\begin{thebibliography}{99}

\bibitem{Jacobson1995}
T.~Jacobson,
``Thermodynamics of spacetime: The Einstein equation of state,''
\textit{Phys.\ Rev.\ Lett.}\ \textbf{75}, 1260 (1995).
\href{https://doi.org/10.1103/PhysRevLett.75.1260}{doi:10.1103/PhysRevLett.75.1260}

\bibitem{MonetteA2026}
K.~Monette,
``Gravitational coupling to entanglement entropy density,''
Independent Research Program, March 2026.

\bibitem{Flammia2012}
S.~T.~Flammia, D.~Gross, Y.-K.~Liu, and J.~Eisert,
``Quantum tomography via compressed sensing: error bounds, sample complexity, and
efficient estimators,''
\textit{New J.\ Phys.}\ \textbf{14}, 095022 (2012).
\href{https://doi.org/10.1088/1367-2630/14/9/095022}{doi:10.1088/1367-2630/14/9/095022}

\bibitem{Paris2009}
M.~G.~A.~Paris,
``Quantum estimation for quantum technology,''
\textit{Int.\ J.\ Quantum Inf.}\ \textbf{7}, 125 (2009).
\href{https://doi.org/10.1142/S0219749909004839}{doi:10.1142/S0219749909004839}

\bibitem{Cotler2020}
J.~Cotler and F.~Wilczek,
``Quantum overlapping tomography,''
\textit{Phys.\ Rev.\ Lett.}\ \textbf{124}, 100401 (2020).
\href{https://doi.org/10.1103/PhysRevLett.124.100401}{doi:10.1103/PhysRevLett.124.100401}

\bibitem{Knill2008}
E.~Knill \textit{et al.},
``Randomized benchmarking of quantum gates,''
\textit{Phys.\ Rev.\ A}\ \textbf{77}, 012307 (2008).

\end{thebibliography}

\end{document}
